% Source: https://artofproblemsolving.com/community/c3h1656197_why_not_100__the_names_bond_james_bond_msa121_msa122
% Stable identity: aops-topic-1656197-q21
Of the eighth graders, $(80-n)\%$ like math. Counting percentages of the whole school,
\[
0.25n+0.35\cdot\dfrac{100+n}{2}+0.40(80-n)=51 .
\]
This is $0.25n+17.5+0.175n+32-0.4n=51$, i.e.\ $0.025n=1.5$, so $n=\boxed{\textbf{(h)}\ 60}$. (Then $60\%$, $80\%$ and $20\%$ of the three grades like math, and $15+28+8=51$.)
